\chapter[Tutorato 2 --- 22 ottobre 2026]{Tutorato 2: estremi di insiemi e verifica di un limite}
\label[appendix]{chap:tutorato-02}

\begin{exercise}[Estremi di un insieme definito da una disequazione]
\label{ex:disequazione-razionale}
Sia
\[
  A:=\left\{x\in\R:\frac{2x-1}{x+2}<1\right\}.
\]
Determinare l'insieme \(A\), i suoi estremi superiore e inferiore
ed eventuali massimo e minimo.
\end{exercise}

\begin{proof}[Soluzione dell'esercizio]
Per \(x\neq-2\), la disequazione equivale a
\[
  \frac{2x-1}{x+2}<1
  \quad\Longleftrightarrow\quad
  \frac{x-3}{x+2}<0.
\]
Il rapporto è negativo se numeratore e denominatore hanno segni
opposti, ossia per \(-2<x<3\). Pertanto
\[
  A=(-2,3).
\]
Per ogni \(\varepsilon>0\), posto
\[
  t:=\min\left\{\frac{\varepsilon}{2},1\right\}>0,
\]
i punti \(-2+t\) e \(3-t\) appartengono ad \(A\) e verificano
\[
  -2<-2+t<-2+\varepsilon,
  \qquad
  3-\varepsilon<3-t<3.
\]
Per i \cref{thm:caratterizzazione-inf,thm:caratterizzazione-sup},
si ha
\[
  \inf A=-2,\qquad \sup A=3.
\]
Poiché \(-2,3\notin A\), per la \cref{def:massimo-minimo}
non esistono né \(\min A\) né \(\max A\).
\end{proof}

\begin{figure}[htbp]
  \centering
  \begin{tikzpicture}[unipd/diagram]
  \draw[->,unipd-muted] (-3.25,0) -- (4.2,0) node[right] {\(\R\)};
  \draw[very thick,unipd-primary] (-2,0) -- (3,0);
  \draw[thick,unipd-primary,fill=white] (-2,0) circle (2.5pt);
  \draw[thick,unipd-primary,fill=white] (3,0) circle (2.5pt);
  \node[below=7pt] at (-2,0) {\(-2\)};
  \node[below=7pt] at (3,0) {\(3\)};
  \node[above=9pt] at (.5,0) {\(A=(-2,3)\)};
\end{tikzpicture}

  \caption{L'insieme \(A=(-2,3)\): gli estremi \(-2\) e \(3\)
  non appartengono all'intervallo.}
  \label{fig:intervallo-aperto}
\end{figure}

\begin{exercise}[Estremi di un insieme definito da termini indicizzati]
\label{ex:estremi-termini}
Sia
\[
  B:=\left\{1-\frac{(-1)^n}{n}:n\in\N,\ n\ge1\right\}.
\]
Determinare gli estremi superiore e inferiore di \(B\) e
stabilire se sono rispettivamente massimo e minimo.
\end{exercise}

\begin{proof}[Soluzione dell'esercizio]
Posto
\[
  b_n:=1-\frac{(-1)^n}{n},\qquad n\ge1,
\]
i primi quattro termini sono
\[
  b_1=2,\quad b_2=\frac12,\quad
  b_3=\frac43,\quad b_4=\frac34.
\]
Si distinguono i due casi, secondo che \(n\) sia pari o dispari:
\[
  b_n=\begin{cases}
    1+\dfrac1n,&n\text{ dispari},\\[4pt]
    1-\dfrac1n,&n\text{ pari}.
  \end{cases}
\]

\emph{Indici dispari.}
Per ogni \(n\ge1\) dispari,
\[
  b_{n+2}=1+\frac1{n+2}<1+\frac1n=b_n.
\]
I termini dispari sono quindi decrescenti e soddisfano
\[
  1<b_n\le b_1=2.
\]

\emph{Indici pari.}
Per ogni \(n\ge2\) pari,
\[
  b_{n+2}=1-\frac1{n+2}>1-\frac1n=b_n.
\]
I termini pari sono quindi crescenti e soddisfano
\[
  \frac12=b_2\le b_n<1.
\]
In entrambi i casi, \(\lvert b_n-1\rvert=1/n\), quindi i termini si
avvicinano a \(1\) al crescere di \(n\). Inoltre,
\[
  \frac12\le b_n\le2\qquad\forall n\ge1.
\]
Poiché \(b_1=2\) e \(b_2=\frac12\), per la
\cref{def:massimo-minimo} si conclude
\[
  \sup B=\max B=2,\qquad
  \inf B=\min B=\frac12.
\]
\end{proof}

\begin{figure}[htbp]
  \centering
  \begin{tikzpicture}[unipd/diagram,x=.75cm,y=2.6cm]
  \draw[->,unipd-muted] (0,.92) -- (10.4,.92) node[right] {\(n\)};
  \draw[->,unipd-muted] (.25,.92) -- (.25,2.18) node[above] {\(b_n\)};
  \draw[dashed,unipd-muted] (.25,1) -- (10,1);
  \node[left=5pt] at (.25,1) {\(1\)};
  \foreach \n in {1,3,5,7,9} {
    \pgfmathsetmacro\v{1+1/\n}
    \fill[unipd-primary] (\n,\v) circle (2pt);
    \draw[unipd-muted] (\n,.9) -- (\n,.94);
    \node[below=5pt] at (\n,.92) {\(\n\)};
  }
  \node[above right=2pt,text=unipd-primary] at (1,2) {\(b_1=2\)};
  \node[above right=2pt,text=unipd-primary] at (3,1.33333) {\(b_3=\frac43\)};
\end{tikzpicture}

  \caption{Termini di indice dispari di \(b_n\), decrescenti
  a partire da \(b_1=2\).}
  \label{fig:termini-dispari}
\end{figure}

\begin{figure}[htbp]
  \centering
  \begin{tikzpicture}[unipd/diagram,x=.75cm,y=3.5cm]
  \draw[->,unipd-muted] (0,.37) -- (11.2,.37) node[right] {\(n\)};
  \draw[->,unipd-muted] (.25,.37) -- (.25,1.13) node[above] {\(b_n\)};
  \draw[dashed,unipd-muted] (.25,1) -- (10.7,1);
  \node[left=5pt] at (.25,1) {\(1\)};
  \foreach \n in {2,4,6,8,10} {
    \pgfmathsetmacro\v{1-1/\n}
    \fill[unipd-primary] (\n,\v) circle (2pt);
    \draw[unipd-muted] (\n,.355) -- (\n,.385);
    \node[below=5pt] at (\n,.37) {\(\n\)};
  }
  \node[above right=2pt,text=unipd-primary] at (2,.5) {\(b_2=\frac12\)};
  \node[below right=2pt,text=unipd-primary] at (4,.75) {\(b_4=\frac34\)};
\end{tikzpicture}

  \caption{Termini di indice pari di \(b_n\), crescenti
  a partire da \(b_2=\frac12\).}
  \label{fig:termini-pari}
\end{figure}

\clearpage

\begin{exercise}[Verifica di un limite finito]
\label{ex:limite-lineare}
Verificare mediante la \cref{def:lez08-limite} che
\[
  \lim_{x\to2}(4x-3)=5.
\]
\end{exercise}

\begin{proof}[Soluzione dell'esercizio]
\emph{Dominio e punto di accumulazione.}
La funzione \(f\colon\R\longrightarrow\R\), definita da
\[
  f(x):=4x-3,
\]
ha dominio \(A=\R\).

Il punto \(2\) è di accumulazione per \(A\). Infatti, per ogni
\(\delta>0\), il numero \(a=2+\delta/2\) appartiene ad \(A\)
e soddisfa
\[
  0<|a-2|=\frac{\delta}{2}<\delta.
\]

\emph{Verifica della definizione.}
Fissato \(\varepsilon>0\), occorre trovare \(\delta_\varepsilon>0\)
tale che
\[
  |f(x)-5|<\varepsilon
  \qquad\forall x\in A:
  0<|x-2|<\delta_\varepsilon.
\]
Per ogni \(x\in A\), si ha
\[
  |f(x)-5|=|4x-8|=4|x-2|.
\]
È pertanto sufficiente scegliere
\[
  \delta_\varepsilon:=\frac{\varepsilon}{4}>0.
\]
Infatti, per ogni \(x\in A\) tale che
\(0<|x-2|<\delta_\varepsilon\), si ottiene
\[
  |f(x)-5|=4|x-2|
  <4\delta_\varepsilon=\varepsilon.
\]
Poiché \(\varepsilon>0\) era arbitrario, per la
\cref{def:lez08-limite} si conclude
\[
  \lim_{x\to2}(4x-3)=5.
\]
\end{proof}
\begin{figure}[htbp]
  \centering
  \begin{tikzpicture}[unipd/diagram,x=2.25cm,y=.57cm]
  \fill[unipd-soft] (1.5,3) rectangle (2.5,7);
  \draw[->,unipd-muted] (.35,0) -- (3.15,0) node[right] {\(x\)};
  \draw[->,unipd-muted] (.5,0) -- (.5,8.6) node[above] {\(y\)};
  \foreach \xx in {1.5,2,2.5} {
    \draw[dashed,unipd-muted] (\xx,0) -- (\xx,7.25);
  }
  \foreach \yy in {3,5,7} {
    \draw[densely dotted,unipd-muted] (.5,\yy) -- (2.82,\yy);
  }
  \draw[very thick,unipd-primary,domain=1.18:2.83,samples=2]
    plot (\x,{4*\x-3});
  \fill[unipd-primary] (2,5) circle (2pt);
  \node[below=4pt] at (1.5,0) {\(2-\delta\)};
  \node[below=4pt] at (2,0) {\(2\)};
  \node[below=4pt] at (2.5,0) {\(2+\delta\)};
  \node[left=5pt] at (.5,3) {\(5-\varepsilon\)};
  \node[left=5pt] at (.5,5) {\(5\)};
  \node[left=5pt] at (.5,7) {\(5+\varepsilon\)};
  \node[above left=3pt,text=unipd-primary] at (2.72,7.88) {\(f(x)=4x-3\)};
\end{tikzpicture}

  \caption{Verifica di \(\lim_{x\to2}(4x-3)=5\):
  per \(\delta=\varepsilon/4\), la condizione
  \(0<|x-2|<\delta\) implica \(|f(x)-5|<\varepsilon\).}
  \label{fig:limite-lineare}
\end{figure}
